%==
%	DOCUMENT DEFINITION
%==

\documentclass[pdftex,11pt,titlepage,twoside,openright]{report}	


% for suporting multi platform
\usepackage[utf8x]{inputenc} 	
\usepackage{amsfonts}
% natbib is what you want for bibliography
\usepackage[square,authoryear]{natbib}


%--
%	DECLARATIONS
%--

%----------------------------------------------------------------------------------------
%	COLOR DEFINITIONS
%----------------------------------------------------------------------------------------

% we use the xtended color package
\usepackage[table,usenames,dvipsnames]{xcolor}

% citation reference color
\definecolor{CITATION_COL}{RGB}{121,0,226}

% urls color
\definecolor{SECTION_COL}{RGB}{0, 117,226}

% internal references color
\definecolor{LINK_COL}{rgb}{0,0,0}

% main color
\definecolor{MAIN_COL}{RGB}{191, 87, 0}
% code highlight color
\definecolor{CODE_COL}{rgb}{0.8,0.8,0.8}

% some black fallback
\definecolor{DEF_COL}{RGB}{226,109,0}

%----------------------------------------------------------------------------------------
%	REFERENCES URL CITATION STYLES AND SO ON
%----------------------------------------------------------------------------------------

% apply color definitions to hyperref package
\usepackage[colorlinks=true,citecolor=CITATION_COL,urlcolor=SECTION_COL,linkcolor=LINK_COL]{hyperref}

% make urls directly open the browser onclick
\usepackage{url}
\urlstyle{same}


%----------------------------------------------------------------------------------------
%	FONT DEFINITIONS
%----------------------------------------------------------------------------------------

% better typography rendering scheme
\usepackage[protrusion=true,expansion=true]{microtype} 

 % we use the standard Times font
\usepackage{pslatex}

 % required for accented characters
\usepackage[T1]{fontenc}

% change line spacing here, readability 
% benefits from a slight increase by default
\linespread{1.25} 

% allows more font size definitions
\usepackage{moresize}

%----------------------------------------------------------------------------------------
%	GEOMETRY  DEFINITIONS
%----------------------------------------------------------------------------------------

% define page styles using geometry
\usepackage[bottom=3cm,marginparwidth=1.5cm, marginparsep=0.5cm]{geometry}

 %define A4 Paper
\geometry{a4paper}	

% define margins for even and odd pages
\setlength{\oddsidemargin}{15.5pt}
\setlength{\evensidemargin}{15.5pt}

% paragraph indent and skip
\setlength{\parindent}{0mm}
\setlength{\parskip}{1mm }

% short command for colorizing a text as a note
% not used within this thesis anymore
\newcommand{\NOTE}[1]{\textit{\textcolor{DEF_COL}{#1}}}
	%\marginpar{\sffamily \MakeUppercase{\tiny{{#1}}}}


%----------------------------------------------------------------------------------------
%	GRAPHICS  DEFINITIONS
%----------------------------------------------------------------------------------------

% required for including pictures
\usepackage[pdftex]{graphicx}

 % make it possible to include more than one captioned figure/table in a single float 
\usepackage{subfig}

% make it possible to control text
% positioning around figures
\usepackage{wrapfig}

 % for much better looking tables
\usepackage{booktabs}	

% for better arrays (eg matrices) in maths		
\usepackage{array} 


%----------------------------------------------------------------------------------------
%	ENVIRONMENT  DEFINITIONS
%----------------------------------------------------------------------------------------

 % very flexible & customisable lists
%(eg. enumerate/itemize, etc.)				
\usepackage{paralist}				

% adds environment for commenting
% out blocks of text & for better verbatim
\usepackage{verbatim}				

% allows multiple rows to be merged to one row
\usepackage{multirow}

%----------------------------------------------------------------------------------------
% CUSTOM STRUT FOR EMPTY BOXES
%----------------------------------------- -----------------------------------------------

% when boxes become a pain this might help
\newcommand{\mystrut}{\rule[-.3\baselineskip]{0pt}{\baselineskip}}

%----------------------------------------------------------------------------------------
%	HEADER / FOOTER  DEFINITIONS
%----------------------------------------------------------------------------------------

% header definitions package
 % this should be set AFTER setting up the page geometry
\usepackage{fancyhdr}

% we want to shorten the display of the chapter in the header
\renewcommand{\chaptermark}[1]
{
  \markboth{#1}{}
}

% but we want to show extended info of the current section
\renewcommand{\sectionmark}[1]{\markright{\thesection\ #1}}

% we can even manipulate the font family
%\lhead{\sffamily Chapter \thechapter}

% this defines, what content should be displayed in header and footer
% and on which side (left-right) as well as on even and odd pages
\fancyhf[FLE,FRO]{}
\fancyhf[HLE,HRO]{\colorbox{MAIN_COL}{\sffamily\LARGE\textcolor{white}{\thepage}}}
\fancyhf[HRE,HLO]{\sffamily \textcolor{MAIN_COL}{\leftmark}}
\cfoot{}

%override plain page style for plain pages
\fancypagestyle{plain}{ %
  \fancyhf{} % remove everything
  \renewcommand{\headrulewidth}{0pt} % remove lines as well
  \renewcommand{\footrulewidth}{0pt}
  \fancyhf[HLE,HRO]{\colorbox{MAIN_COL}{\sffamily\large\textcolor{white}{\thepage}}}
  \cfoot{}
}

%----------------------------------------------------------------------------------------
%	SECTION TITLE APPEARANCE
%----------------------------------------------------------------------------------------

% we want to customize our section title 
% to make it look more appealing
\usepackage{sectsty}

% for dimensioning and other math stuff
\usepackage{amsmath}

% modify bg colors
\usepackage[some]{background}

% draw custom graphics progammatically
\usepackage{tikz}				
\usetikzlibrary{shapes, backgrounds,mindmap, trees}

% this is our custom section page layout
\newcommand{\SECTBOX}[1]
{
	\begin{tikzpicture}[remember picture,overlay]
 		\path [fill=MAIN_COL] (current page.west)rectangle (current page.north east); 
	\end{tikzpicture}
}

% this is our custom title page layout
\newcommand{\TITLEBOX}
{
	\begin{tikzpicture}[remember picture,overlay]
 		\path [fill=MAIN_COL] (current page.north west) rectangle (25.0,-3.5); 
	\end{tikzpicture}
}


 % (See the fntguide.pdf for font help)
\allsectionsfont{\sffamily\mdseries\upshape}

% for customizing the title page
\usepackage[explicit]{titlesec}

% font and color def for section
\titleformat{\section}{\LARGE\sffamily}{\textcolor{LINK_COL}{\thesection \hspace{12pt} #1}}{0pt}{}

% font and color def for subsection
\titleformat{\subsection}{\large\bf\sffamily}{\textcolor{SECTION_COL}{\thesubsection \hspace{12pt} #1}}{0pt}{}

% note: remove \thesection or \thesubsection to leave out the 
% numbering at the start of the section / subsection headline


%----------------------------------------------------------------------------------------
%	TABLE OF CONTENT APPEARANCE
%----------------------------------------------------------------------------------------

% deprecated
%\usepackage[nottoc,notlof,notlot]{tocbibind} % Put the bibliography in the ToC

% Alter the style of the Table of Contents
%\usepackage[titles,subfigure]{tocloft}

%\renewcommand{\cftsecfont}{\rmfamily\mdseries\upshape}

%\renewcommand{\cftsecpagefont}{\rmfamily\mdseries\upshape} % No bold!

% renames abstract to executive summary
\renewcommand{\abstractname}{Executive Summary} 		


%----------------------------------------------------------------------------------------
%	CUSTOM HRULE
%----------------------------------------------------------------------------------------

% a custom horizontal rule / line
\newcommand{\HRule}{\rule{\linewidth}{0.5mm}}

\newcommand{\ThinHRule}{\textcolor{CODE_COL}{\rule{0.5\linewidth}{0.05mm}}}


%----------------------------------------------------------------------------------------
%	ENVIRONMENTS
%----------------------------------------------------------------------------------------

% extended listin options
% can be used for code display
\usepackage{listings}

% syntax highlight def for listings
\lstset{ %
  	backgroundcolor=\color{CODE_COL},
	language=C,
	breaklines=true
}

% allows us to define pseudocode algorithms
\usepackage[ruled,linesnumbered]{algorithm2e}

% extended table definitions
\usepackage{tabularx}

% include external pdf pages
\usepackage{pdfpages}
\usepackage{float}


%----------------------------------------------------------------------------------------
%	Code Syntax Highlight Formatting
%----------------------------------------------------------------------------------------

\lstset{
  basicstyle=\ttfamily,
  columns=fullflexible,
  showstringspaces=false,
  commentstyle=\color{gray}\upshape
}

\lstdefinelanguage{XML}
{
  morestring=[b]",
  morestring=[s]{>}{<},
  morecomment= [s]{<!}{->},
  stringstyle=\color{black},
  identifierstyle=\color{blue},
  keywordstyle=\color{orange},
  morekeywords={xmlns, pattern,type, cells, count, required} % list your attributes here which may be highlighted
}






%--
%	INDEX AND GLOSSARY
%--


\usepackage[style=long,nonumberlist,toc,xindy,acronym,nomain]{glossaries} % nomain, if you define glossaries in a file, and you use \include{INP-00-glossary}
%\loadglsentries[main]{glossary}
% or using \input:
\include{glossary}
\makeglossaries
\usepackage{makeidx}
\makeindex



%==
%	BEGIN DOCUMENT
%==


\begin{document}

% roman numbering 
\setcounter{page}{1}
\pagenumbering{roman}

% print title
\begin{titlepage}


%TIKZ BACKGROUND
\TITLEBOX

%CONTENT
\begin{center}



\sffamily\textsc{\huge{\textcolor{white}{The University of Texas at Austin}}}\\[4cm]


% Upper part of the page. The '~' is needed because \\
% only works if a paragraph has started.
\includegraphics[width=0.35\textwidth]{media/faculty.png}\makebox[1.5cm]{}\includegraphics[width=0.35\textwidth]{media/unilogo.png}~\\[1.5cm]



% Title
\HRule \\[0.4cm]
{
 \huge \bfseries \sffamily Vision Transformer-Assisted Analysis of
 
Neural Image Compression and Generation  \\[0.4cm] 
}

\HRule \\[0.4cm] 

\normalfont \LARGE \sffamily Master's Thesis Report \\[1.5cm]
% \normalfont \large \sffamily DRAFT v2.0 \\[0.5cm]
\normalfont \large \sffamily Official Code Repositories: \\
\normalfont \large \sffamily \url{https://github.com/kliment-slice/thesis-code} \\[0.05cm]
\normalfont \large \sffamily \url{https://github.com/kliment-slice/thesis-latex} \\[1.0cm]

% Author and supervisor
\noindent
% \begin{minipage}{0.4\textwidth}
% \begin{flushleft} \large
% \emph{Supervisors:} \\
% Prof. Dr. Alan C. Bovik\\
% Prof. Dr. Rachel A. Ward\\
% \end{flushleft}
% \end{minipage}%
\begin{minipage}{0.4\textwidth}
\begin{center} \large
\emph{Author:}\\
Kliment Minchev\\
\end{center}
\end{minipage}

\vfill


% Bottom of the page
{\large \today}

\vfill



% \vfill
\end{center}
\end{titlepage}



%pagestyle after title
\pagestyle{fancy}

% table of contents,list of figures and list of tables
\setcounter{tocdepth}{1}
% \tableofcontents
% \listoffigures
% \listoftables
{\tableofcontents \let\cleardoublepage\clearpage \listoffigures 
\let\cleardoublepage\clearpage \listoftables \let\cleardoublepage\clearpage}

%\cleardoublepage

% ABSTRACT EXECUTIVE SUMMARY

\begin{abstract}
    This work investigates a novel application of a Vision Transformer (ViT) as a 
    quality assessment reference metric for generated images after neural image compression. 
    The Vision Transformer is a revolutionary implementation of the Transformer attention 
    mechanism (typically used in language models) to object detection in digital images. 
    The ViT architecture is designed to output a classification probability distribution against
    a set of training labels. 
    Thus, it can be a suitable candidate for a new method for quantitative assessment of generated 
    images based on object-level deviations from the original pre-compression image. 
    This approach complements other comparative measurement techniques based on per-pixel 
    discrepancies (Mean Squared Error, MSE) or structural comparison (Structural Similarity Index, SSIM). 
    This study proposes an original end-to-end deep learning framework for neural image compression, 
    latent vector representation, regeneration, and image quality analysis using 
    state-of-the-art model architectures. Neural image compression and generation 
    is achieved using a Generative Adversarial Network (GAN). 
    Results from this work demonstrate that a Vision Transformer is capable of assessing 
    the quality of a neurally compressed image. 
    Moreover, this methodology provides valuable insights when measuring image quality.
    It can be used in addition to established perceived quality metrics for compressed and generated 
    images such as SSIM and Frechet Inception Distance (FID). 
    
    \ThinHRule
\end{abstract}

\pagestyle{fancy}

\setcounter{page}{1}
\pagenumbering{arabic}


%==
%	CHAPTER 1
%==

\chapter{Introduction to Vision Transformers (ViT)}


% Chapter overview 
This chapter presents the reader with an introduction to Vision Transformers (ViT).

It covers the motivation as to why ViT, or a future development inspired by it,
will have a profound impact on image compression, analysis, and generation overall.
This chapter presents evidence that a Transformer, or perhaps an evolved deep learning model with a
similar architecture, i.e. generalizable and highly overparameterized, can be superior in
compressing and evaluating the latent feature space of a digital image compared to present-day
technologies.

This section then reviews a brief history or Transformer usage in deep learning.
These generalized architectures are dominating state-of-the-art language models,
as they are extremely efficient in packing information within a one dimensional vector.

This introduction will then proceed to describe the principles of operation of a ViT. Then, proceed with
a mathematical formulation. Finally, it will cover currently available implementations
in the form of pre-trained models and conclude with an explanation on the computational and
financial constraints of training such demanding architectures.

\ThinHRule

\newpage
%--
%	CHAPTER 1
%--

\section{Motivation}

Transformers are presently considered to hold a great promise for the future of Deep Learning
as a step towards Artificial General Intelligence.
Due to their architecture, they are more generalizable, less prone to overfitting, and able 
to learn highly complex representations. The Transformer architecture has already been proven 
to make obsolete Recurrent Neural Networks (RNNs) in natural language models. Furthermore, the 
Vision Transformer (ViT) has outperformed certain Convolutional Neural Networks (CNNs) in image classification tasks. \citep{dosovitskiy2020vit}

Figure 1.1 below shows an increase in the popularity of research related to Vision Transformers.

\begin{figure}[H]
	\begin{center}
	\includegraphics[width=1\textwidth]{media/papersWcode_tfUsage2.png}
	\end{center}
	\caption[Historical Usage of ViT in Image Tasks]{As of 2022, the usage of a Vision Transformer (ViT) in image 
	tasks matches the usage of ResNets and has outnumbered any other popular CNN architecture.
	\citep{PapersOverTime}}
	\end{figure}

Figure 1.1 was produced by PapersWithCode, a popular academic research aggregator. For the past three years, 
ResNets, the most popular architecture in image processing and computer vision, has dominated the proportion
of academic research in object detection. In 2022, Vision Transformer research popularity has reached
that of ResNets and exceeded any other major category.


In the zeitgeist of Vision Transformer research, this thesis explores a ViT-assisted metric related to 
image compression. This metric can provide additional insights to GAN output quality 
and the latent space (contextual) preservation of a variety of input images.


\NOTE{Thus, contributions from this thesis can be viewed as providing a stepping stone 
towards an end-to-end Transformer-based image compression and reconstruction framework.}


\section{Brief History}

\subsection{Attention and Language Models}

"Attention Is All You Need" is a seminal research publication
by a team of Google researchers, which kickstarted the Transformer revolution in Deep Learning in 2017.
It proposes a novel architecture, which models long-range dependencies in
sequential (text) data, by arranging a set of self-attention layers. 

A self-attention layer is what the model uses to focus on different elements of the 
input sequence simultaneously. For example, it can be used to compute the 
distance (relationship) between every word in a given sentence. \citep{Attention}


Examples of implementations of text-based Transformers are BERT by Google and GPT-3 

by OpenAI.
BERT, among many other applications, processes and autofills every single English-based 
Google user search query as of 2021 . \citep{bert}
GPT-3, on the other hand, revolutionized text generation in 2020, demonstrating the ability 
to generate extremely cohesive textual output.

Most Transformers are used for applications in language modeling and Natural Language 

Processing (NLP).
Thus, they are often benchmarked against Recurrent Neural Networks (RNNs, and specifically 
Long Short-Term Memory, LSTM architecture). LSTMs rely on hidden states to pass information 
along sequentially during the encoding and decoding process for each word token. 
However, they typically fall short learning long-range dependencies.


\subsection{Attention in Vision Tasks}

The attention mechanism is capable of focusing on objects found anywhere on an input image.
It operates within a single network layer compared to Convolutional Neural Networks (CNNs),
where the variable size convolution kernels scan across the different layers of the architecture.
\citep{dosovitskiy2020vit}

As shown in Figure 1.2 below, tokenization happens at the pixel level, i.e. each pixel attends to each other pixel in the grid. 
This becomes computationally intensive, on the order of $(n^2)^2$, where $n$ denotes width of a square image. 
To resolve this, the input image is broken down into square blocks of equal size, referred to as image patches.
Then, each image patch is unrolled into a one-dimensional sequence $(n x 1)$ and indexed with a positional 
embedding in a table for future reference and retrieval purposes. The embeddings enter the Transformer and 
finally, a feed forward classifier, in the form of a Multilayer Perceptron (MLP) 
makes the classification prediction, yielding a probability distribution. \citep{dosovitskiy2020vit}

\begin{figure}[H]
	\begin{center}
	\includegraphics[width=0.8\textwidth]{media/vit_arch.png}
	\end{center}
	\caption[ViT Architecture]{The Vision Transformer (ViT) architecture. \citep{GoogleViT}}
	\end{figure}

A Transformer, in a way, is a generalization of a feed forward network, but 
instead of fixed connections weights in an MLP, each connection weight (i.e. attention) 
is computed ad hoc. This makes the Transformer, unlike the MLP, permutation invariant. 
That is, it would not know where certain information is coming from, unless there are additional 
learnable sequential positional embeddings, i.e. index the image patches.

\section{Principles of Operation}

Continuing from the previous section, a ViT can be thought of as a 
generalization of an MLP, which itself is a generalization of a CNN. 
The ViT happens to learn very similarly to a CNN, which represents the latent space
as filters carrying principal components.

In principle, CNNs have good inductive priors and can learn any function. 
However, they promote locality, i.e. nearing pixels are probability-wise considered most important.
This may easily not be desired, especially in the key applications of object detection and, in the 
future, image compression.


The encoding process indexes embeddings. For instance, certain key words in a sentence or 
objects in image blocks are mapped in a reference lookup table. 
The Decoder outputs Keys at each step. These vectors represent hidden states,
which are being passed on into each next iteration of the Transformer. 
The last layer, expectedly, uses a Softmax architecture to normalize and map the potential output classes 
to a probability distribution. 

%\vspace{10mm}
\newpage

\textbf{Multi-Head Attention}


As shown in Figure 1.3 below, sets of parallel attention layers at each token are called multi-head attention.
This approach varies what to pay attention to, for example,at the different objects 
in an image (or in natural langue, different verbs in a sentence). 
The multi-head attention is composed of Key-Value pairs coming from 
the encoding part of the source image (i.e. the input embedding) and 
Queries from the output embedding (i.e. encoding part of target image).

\begin{figure}[H]
	\begin{center}
	\includegraphics[width=0.5\textwidth]{media/attn_diag.png}
	\end{center}
	\caption[Attention Mechanism]{The Attention mechanism architecture. \citep{GoogleViT}}
	\end{figure}


\section{Mathematical Formulation}

%\textbf{Attention}
\subsection{Attention}

In its full formulation, Attention is a function of vectors representing Queries, Keys, and Values, 
labeled as $(Q,K,V)$. 
Attention equals the dot product $(QK^T)$ of Keys and Queries respectively, softmaxed over 
the square root of dimensions and multiplied by the Values vector.

To provide further intuition:

Values are what is most interesting in the source image, i.e. 
attributes or structural features. In text, a Value could be important adjectives before each keyword, 
which provide emphasis in a given input sentence.
Keys, on the other hand, index (or address) those Values (e.g. name, type, weight). 
Each Key has an associated Value. 
Queries are built by the encoder of the target image and prompt the network to find closest 
available information (Key and its corresponding Value). 

\begin{center}
	$ Attention(Q, K, V) = softmax(\dfrac{QK^T}{\sqrt{d_k}} V ) $
	\end{center}


\NOTE{Thus, the overall dynamic is that a Query is pegged against a Key to locate a certain Value.}

%\vspace{10mm}

%\textbf{Softmax}
\subsection{Softmax}

The Softmax function is defined as a normalized exponential function. 
A sequence of variables is mapped into exponentials and divided by the sum of all exponentials. 
Thus, the large numbers become almost ones and small numbers near zeros. Softmax is similar to the maximum 
function, with a key difference that Softmax is differentiable. 

\begin{center}
	$ \displaystyle\sigma(Z_i) = \dfrac{e^{z_j}}{\sum_{j=1}^K e^{z_j}}$ for $i=1,...,K$ $and z=(z_1,...,z_K) \in \mathbb{R}^K $
	\end{center}

Thus, a Softmax of an inner product of each Key with Query vector normalizes to a 
probability distribution over all Values. Neural networks typically utilize a softmax in their last 
layer over all the classification labels. This yields the top classification by probability. 
Using the softmax function, a certain Key will stand out (in magnitude) vs the rest.

%\vspace{10mm}

%\textbf{Vector Similarity}
\subsection{Vector Similarity}

Vector proximity between embeddings represents the similarity between objects in images, or 
word connotations in sentences.


\begin{figure}[H]
	\begin{center}
	\includegraphics[width=0.5\textwidth]{media/dot_real.png}
	\end{center}
	\caption[Vector Representation]{Vector representation on a unit circle.}
	\end{figure}


The dot product of Keys and Queries yields an angle between both vectors (e.g. $u_i$ and $u_j$ in Figure 1.2 above) to measure 
how similarly aligned they are. In high dimensional spaces, most vectors would be orthonormal to each other
and $cos(90)=0$. But if Key and Query vectors are similar or align, they'd have a large dot product. 
The larger the similarity, the larger the dot product. 
The Query vector is computed with each Key in the surrounding vector space and softmaxed to 
select the one Key with the highest dot product.
The selected Key in space has an associated Value. In applications, that Value would correspond to 
a match with a labeled object in an image, or perhaps the next word in a generated sentence. 


A depiction of the vector space is shown in Figure 1.3 below.


\begin{figure}[H]
	\begin{center}
	\includegraphics[width=0.3\textwidth]{media/dot2_real.png}
	\end{center}
	\caption[Key/Query Vector Proximity]{Vector proximity shows the closest Key vector to a given Query vector in space.}
	\end{figure}


\section{Implementations}

This section reviews several notable ViT implementations of interest.
Most developments have been made open source. Howerver, some of the 
highest quality implementations are still closed source and operated 
under a license or payment wall.

\subsection{Open Source}

The academic research aggregator PapersWithCode lists 96 open source implementations of 
the original ViT from "An Image is Worth 16x16 Words" discussed in Chapter 2. 
The publication was made for ICLR 2021. \citep{Papers16}

The original model from the team of Google researchers, written using TensorFlow, has 72,567 stars in its GitHub repository.
The second highest implementation by Hugging Face, written using PyTorch, has 61,820 stars. 


\vspace{5mm}

\textbf{Pre-trained Model Used}

For the purpose of this thesis, a PyTorch implementation was used trained on ImageNet-21k
and fine tuned on ImageNet-1k. The ViT has 9,696 stars on its GitHub repository,
ranking as the 5th highest rated implementation. It was chosen due to its reliability,
project maturity, and interfaceability. \citep{Papers16}

\vspace{5mm}

A python pip package is also available on PyPI from the same "lucidrains" implementation,
which can be installed via: 
\begin{verbatim} "pip install pytorch_pretrained_vit" \end{verbatim}


Other pre-trained models are also availabe on GitHub and Google Colab. 
Some suffer from a lower GitHub star rating, are not well packaged (not object-oriented), or perhaps deviate too far
from the original implementation of "An Image is Worth 16x16 Words". 


\subsection{Closed Source}


Several modified implementations targeting specific generative applications have been made by Hugging Face and OpenAI.
OpenAI has expanded from its success from GPT-3, a text generation Transformer, to visual applications.
Transformers and Vision Transformers can be used to find context in image objects from the latent space of an input image.
CLIP, Image GPT, DALL.E, and DALL.E2 by OpenAI bridge the gap between textual information and images.
DALL.E 2, the highest quality image Transformer to date, was released in April 2022.
It is trained on 250M image-text pairs to be able to generate coherent images from textual description. \citep{Dalle2}

\NOTE{Most of OpenAI's Transformer developments (e.g. GPT-3, DALL.E 2) have not been open sourced. 
Google's ViT and BERT, however, have been.}


\section{Computational Constraints for Training}


The main reason for using a pre-trained ViT is that the Transformer architecture is so generalizable and 
overparameterized, that the computational requirements for Time, Memory, and Cost of training 
are beyond the scope of this thesis. 


For example, GPT-3 was trained on 45TB of text data (e.g. all of Wikipedia included). It has 175B parameters and 96 attention layers. \citep{GPT3}

One flavor of GPT-3 (of many) alone cost \$4.6M to train using the most advanced Nvidia datacenter GPUs in a cluster.

The original ViT implementation from "An Image is Worth 16x16 Words" took 2.5k core-days of training time. \citep{dosovitskiy2020vit} 

\vspace{5mm}

\textbf{Training the Next Generation of Image GPT}

It would require several orders of magnitude more resources to train a Vision Transformer that matches the quality of GPT-3 
for generative text applications.

The training data would need to be similar to all Google images on the internet (or at least a respectable representative subset).
It would probably cost on the order of \$50M in training alone, since image matrices are two-dimensional sequences of data.
Contextual complexity also rises exponentially for images compared to text.













%==
%	CHAPTER 2
%==
% \cleardoublepage
\chapter{Literature Review: Transformers and Neural Image Compression}


This chapter serves to present the audience with a literature review of seminal academic
publications relating Transformers to image compression and generation.

Other relevant background information is also provided to the audience
to aid the reader's understanding of this thesis project.

Among much of the domain knowledge available, the reader would find interest in the 
takeaways offered from:

\begin{itemize}
	\item "An Image is Worth 16x16 Words" \citep{dosovitskiy2020vit}
	\item "Towards End-to-End Image Compression and Analysis with Transformers" 
    
    \citep{Bai2022AAAI}
    \item Image Generation
	\item "TransGAN: Two Pure Transformers Can Make One Strong GAN" \citep{jiang2021transgan}
	\item "First Principles of Deep Learning and Compression" \citep{Principles}
	\item Review of commonly used Image Quality Assessment (IQA) metrics
    
    \citep{Metrics}
\end{itemize}

\ThinHRule

\newpage
\section{"An Image is Worth 16x16 Words"}

\large{Transformers for Image Recognition at Scale}

Now that we explained relevant transformer components, we can see how it applies 
to 2D signals, i.e. image matrices for classification purposes in image recognition.
The Vision Transformer or ViT, was very recently published for ICLR 2021 by a 
google team for a ImageNet trained transformer.

The attention mechanism is capable of focusing on objects anywhere on the image 
within a single network layer versus the variable size convolution kernels 
across the different layers.

"image is worth 16x16 words" completely discards the notion of convolutions. 
Compared to a convolutional counterpart (the ResNet), the ViT cost 75\% 
less to train and beats accuracy on ImageNet by 1\%. 
ViT uses approximately 2-4x less compute to attain the same performance 
(averaged over 5 datasets). Unlike variable size convolution kernels accross 
layers, the block size in the image transformer is able to pay attention within 
a single layer to anywhere on the image. 
CNNs have good inductive priors and can learn any function. 
However, this promotes locality (i.e. nearing pixels are probably most important).
One way to think of a ViT is a generalization of an MLP which itself is a 
generalization of a CNN. The ViT also learns very similarly to a CNN 
(e.g. filters with principal components).
Transformer, in a way, is a generalization of a feed forward network, but 
instead of fixed connections weights in an MLP, each connection weight (i.e. attention) 
is computed on the fly. That makes the Transformer, unlike the MLP, permutation invariant. 
In that it wouldn't know where information is coming from unless there are additional 
learnable sequential positional embeddings, i.e. number down the image patches.

% \begin{figure}[H]
% \begin{algorithm}[H]
% \KwData{A as Array to sort,}
% \KwResult{A as sorted Array} 

% int n $\leftarrow$ A.size \tcp*[l]{cache the initial size of A}


% \Repeat{n>1}{
% 	int newn $\leftarrow$ 1
% 	\For{int i=0; i<n-1; i++}
% 	{
% 		\If{A[i] > A[i+1]}
% 		{
% 			A.swap(i, +1);
% 		}
% 	}
% 	n  $\leftarrow$ newn
% }

% return A
% \caption{bubbleSort(Array A)}
% \end{algorithm}
%\caption[TBC]{Figure caption.v.}
% \end{figure}


\newpage
\section{"Towards End-to-End Image Compression with Transformers"}


% \begin{table}[H]
% \begin{center}
% \begin{tabular}{|c|c|c|c|c|c|c|}
% \hline
% result XY & \multicolumn{6}{c|}{Rater A}\\\hline
%  \multirow{6}{*}{Rater B}& & \textbf{9} & \textbf{0} & \textbf{1} & \textbf{2} & \textbf{c(a) }\\\cline{2-7}
%  & \textbf{9}& 27	&0&	0	&0 & 0.5510204082\\\cline{2-7}
% &\textbf{0}	&0	&6	&0	&0&0.1224489796\\\cline{2-7}
% &\textbf{1}	&0	&1	&7	&0&0.1632653061\\\cline{2-7}
% &\textbf{2}	&0	&1	&7	&0&0.1632653061\\\cline{2-7}
% &\textbf{c(r)} & 0.5510204082 & 0.1632653061 & 0.2857142857 & 0& n=49\\\hline

% \end{tabular}
% \caption[Example data for mathematical equations]{A set of example data for further use in a mathematical equation.  }
% \end{center}
% \end{table}


\section{Image Generation with GANs}

Generative models.
GANs used to be the coolest thing until Transformers came along, but nonetheless 
still great. Typically, GANs pit 2 CNNs against each other, one called Generator to 
generate an image and one called Discriminator to determine how fake the image is 
compared to the input, relaying feedback to improve the Generator after each epoch. 
A GAN pits two Neural Networks (e.g. Convolutional) one against the other to generate 
new content (e.g. an image). The structure of a GAN is foundational to its design. 
As Fig.1 shows, the Generator CNN is shown on the left of the generated output. 
The Generator iteratively learns to create new synthetic data resembling real source data 
(e.g. regular or 360° images).  On the right of the generated output is the Discriminator CNN.
The Discriminator is trained to differentiate between real and synthetic data. 
In essence, the Discriminator gives the Generator feedback on its performance while 
simultaneously achieving incremental improvement in discernibility. 
Thus, each epoch the Generator produces more and more realistic data and the Discriminator iteratively improves its own ability to differentiate between real and synthetic data.


\section{First Principles of Neural Image Compression}

Deterministic, Probabilistic
GANs
JPEG/MPEG

\section{Commonly Used Metrics for Image Quality}
Metrics for the generated output image are: a BRISQUE score (reference-less) [5], 
Mean Squared Error (MSE) and Structural Similarity Index (SSIM). PSNR.
GANs use FID and Inception score. 
BRISQUE \citep{BRISQUE}
Divisive normalization, Gaussian properties from natural scene statistics
Process results in Gaussian unless it's distorted
BRISQUE too simple for user-generated content (UGC)

%==
%	CHAPTER 3: Core of Thesis
%==
% \cleardoublepage
\chapter{ViT-based Assessment of Neural Image Compression}

Now that we explained relevant transformer components, we can see how it applies to 
2D signals, i.e. image matrices for classification purposes in image recognition.
The Vision Transformer or ViT, was very recently published for ICLR 2021 by a google 
team for a ImageNet trained transformer.
Tokenization happens at pixel level, so each pixel would have to attend to each other 
pixel in the grid, which becomes too heavy to compute, on the order of. 
To resolve that, the image is broken down into blocks of equal size, a 16x16 subset 
of the image called image patches. Then, unroll each image patch into a sequence (256x1), 
and index it with a positional embedding in a table. All of that is then fed into a 
standard Transformer, like from Attention is all you need. Finally, a feed forward 
classifier (MLP) makes the classification prediction, voila image recognition.
The total number of parameters is on the order of 100M.

\ThinHRule

\newpage
\section{Generative Image Compression and Generation}

(vineeth)

\subsection{Architecture}

SGD optimizer for GAN

\subsection{Sample Images Used}

\begin{figure}[H]
	\begin{center}
	\includegraphics[width=0.7\textwidth]{media/input_images_labeled.png}
	\end{center}
	\caption[Input Images]{Input Images used in this project (512x512):
    \newline
    \hspace{10pt} "Bevo": The University of Texas mascot, a famous longhorn bull.\newline
    \hspace{10pt} "Kliment": A face portrait of the author. \newline
    \hspace{10pt} "Logo": The Texas Longhorns logo.\newline
    \hspace{10pt} "Tower": The University of Texas Tower, the Main Building on campus.}
	\end{figure}

        % \begin{figure}[H]
% \begin{algorithm}[H]
% \KwData{A as Array to sort,}
% \KwResult{A as sorted Array} 

% int n $\leftarrow$ A.size \tcp*[l]{cache the initial size of A}


% \Repeat{n>1}{
% 	int newn $\leftarrow$ 1
% 	\For{int i=0; i<n-1; i++}
% 	{
% 		\If{A[i] > A[i+1]}
% 		{
% 			A.swap(i, +1);
% 		}
% 	}
% 	n  $\leftarrow$ newn
% }

% return A
% \caption{Neural Image Compression}
% \end{algorithm}
% \caption[TBC]{Figure caption caption}
% \end{figure}

\subsection{Latent space vector representation during compression}

$nx1$ vector, where $n$ corresponds to the height or width (in pixels)
of a square input image. In the case of all input images used, the latent vector
is of size $512x1$, since the input images are of size $512x512$.

hard to decipher but here is a zoomed in visual of the first
This is what the GAN architecture compresses the full image to (1.54kb vs original size 409kb).
The GAN then rebuils to 242kb.


\begin{figure}[H]
	\begin{center}
	\includegraphics[width=0.5\textwidth]{media/latent.png}
	\end{center}
	\caption[Latent Space Representation]{Visual representation of the first 10 pixels in the most compressed version of the original image.}
	\end{figure}

% \newpage
\section{Output and Visual Inspection}
The generative process from the GAN was designed to output an image at a specified epoch.

\subsection{Training Process}
Below is a demonstration of the GAN learning process at every 250 epochs:
\begin{figure}[H]
	\begin{center}
	\includegraphics[width=0.8\textwidth]{media/ganLearn.png}
	\end{center}
	\caption[GAN Training Process]{The GAN learns to compress and generate the Texas Longhorns logo.}
	\end{figure}

Due to the cutting edge nature of the technologies used, a natural performance asymptote
was observed. The GAN was able to reconstruct certain input images better than others.
After a certain iteration, as usual, the GAN was unable to further learn how to compress,
represent, and regenerate some input images. Typically, once a GAN reaches this stage,
it learns from random noise and generation performance decreases.

\subsection{Generated Results}

Figure 3.4 below shows the resulting images from the neural compression GAN.

\begin{figure}[H]
	\begin{center}
	\includegraphics[width=0.8\textwidth]{media/gan_compress2.png}
	\end{center}
	\caption[Neural Compression and Generation]{After sufficient training, the GAN outputs a regenerated
    version of the original image from a latent space vector representation.}
	\end{figure}

The GAN was able to respectably regenerate "Kliment" and "Logo", especially if the resolution
were to be lowered (e.g. 32x32, such as the CIFAR-10 dataset).

However, the GAN was unable to perform as well for "Bevo" and "Tower".
It learned random noise and generation performance decreased.

\section{ViT-Scores}

The Vision Transformer-Assisted ViT Score is an original development from this thesis. 

It is an attempt to measure the quality of a generated image after neural compression.

The ViT-score is in the open interval $(0,1)$ with $0$ being poor and extremely dissimilar from
the original and $1$ being excellent and fully similar to original.

Mathematically, the endpoint values of the interval are unattainable by probabilistic models
such as the GAN.

\subsection{Mathematical Formulation}

The following is a mathematical representation explained in further detail.

\begin{center}
    $ViT_{score} = \displaystyle\dfrac{argmax_{A'\subset A,\lvert A' \rvert = k } \sum_{a \in A'} {a} }{k} $
\end{center}


where $\sum_{a \in A'} {a}  = \lbrace{m \in I_{input}}\rbrace \cap \lbrace{n \in I_{generated}}\rbrace$


and $m$ are the top-$K$ labels in the input image $I_{input}$

and $n$ are the top-$K$ labels in the generated image $I_{generated}$.

\vspace{12pt}

This overly elaborate mathematical notation is an attempt at describing:


\NOTE{
"From the full set of trained ViT labels, we find the top-$K$ number of intersecting labels between the original and generated images. Then, we divide that by $K$"
}


For example, of the top-100 labels found in the original image, 
identify the set of labels also found in the generated image. 
Then, divide that number of intersecting
labels by the total number of 100 labels.


\subsection{ViT-Scores from Resulting Images}


Following Figure 3.7, the ViT-scores for the GAN generated images after neural compression
are as follows:

\begin{center}
	\includegraphics[width=0.6\textwidth]{media/gan_compress2.png}
\end{center}

\begin{table}[H]
\begin{center}
\begin{tabular}{|c|c|}
\hline
\textbf{Bevo}	& 0.14\\
\textbf{Kliment}	& 0.54\\
\textbf{Logo}	& 0.29\\
\textbf{Tower} & 0.03\\\hline

\end{tabular}
\caption[ViT-Scores of Generated Images]{ViT-Scores demonstrate a somewhat expected quality assessment.}
\end{center}
\end{table}



\textbf{"Kliment"} leads with a ViT-score of $0.54$, which is understandable as the GAN 
generated a face (although smudgy) and was rather able to recreate the scenery structurally.

\textbf{"Logo"} generation seems structurally excellent and the ViT-score is $0.29$, which
is considered a good score for this particular GAN architecture and training.

\textbf{"Bevo"} barely preserves the original shape at ViT-score of $0.14$, while
\textbf{"Tower"} is incomprehensible and barely resembles the original at ViT-score of $0.03$.


Overall, the ViT-score does a good job of measuring image quality.


\section{Established IQA Metrics}

\subsection{"Kliment"}

\begin{figure}[H]
	\begin{center}
	\includegraphics[width=1\textwidth]{media/kimbo_metrics.png}
	\end{center}
	\caption["Kliment" Established Metrics]{After sufficient training, the GAN outputs a regenerated
    version of the original image from a latent space vector representation.}
	\end{figure}

\subsection{"Logo"}

\begin{figure}[H]
	\begin{center}
	\includegraphics[width=1\textwidth]{media/logo_metrics.png}
	\end{center}
	\caption["Logo" Established Metrics]{After sufficient training, the GAN outputs a regenerated
    version of the original image from a latent space vector representation.}
	\end{figure}


\subsection{"Bevo"}

\begin{figure}[H]
        \begin{center}
        \includegraphics[width=1\textwidth]{media/bevo_metrics.png}
        \end{center}
        \caption["Bevo" Established Metrics]{After sufficient training, the GAN outputs a regenerated
        version of the original image from a latent space vector representation.}
        \end{figure}

\subsection{"Tower"}

\begin{figure}[H]
        \begin{center}
        \includegraphics[width=1\textwidth]{media/tower_metrics.png}
        \end{center}
        \caption["Tower" Established Metrics]{After sufficient training, the GAN outputs a regenerated
        version of the original image from a latent space vector representation.}
        \end{figure}

\subsection{"BRISQUE"}

Blind/Referenceless Image Spatial Quality Evaluator (BRISQUE)
where approaching 0 is a good score and approaching 100 is a bad score, the BRISQUE referenceless 
image quality methodology

This score could be interpreted as the image being more photorealistic than not. 
In terms of quality, this compares to a camera captured image with quality corruption 
caused by blurs or graininess. An image with no distortions often has a score below 5.

\begin{table}[H]
    \begin{center}
    \begin{tabular}{|c|c|c|}
    \hline
      & Original & Generated \\ [0.5ex] 
    \hline\hline
    \textbf{Bevo}	& 32.9214 & 39.5535\\
    \textbf{Kliment}	& -8.3593 & 44.3570\\
    \textbf{Logo}	& 102.9010 & 97.1844\\
    \textbf{Tower} & 14.5973 & 52.8363\\\hline
    
    \end{tabular}
    \caption[BRISQUE]{BRISQUE Scores of original and generated images.}
    \end{center}
    \end{table}

Expectedly, the BRISQUE values for the generated images are always higher than
their original counterparts. "Logo" is not a photorealistic image to begin with, 
so it is understandable that the BRISQUE value is high at $102.9$.
None of the generated images would pass BRISQUE as photorealistic and free of 
distortions.  

Loss functions as well (MSE loss was used in GAN)


\section{GAN-Related Quantitative Metrics}


(FID score, inception score (IS) )
Can be loss functions as well (MSE was used)
FID is Frechet Inception Distance. 0 if there is no difference between the images.
81105.162 for logo and logo\_GAN.
331171.556 for tower and towerGAN
549089.491 for kimbo and kimboGAN
1241999.901 for bevo and bevoGAN


\section{Summary of Results}

\begin{center}
	\includegraphics[width=1\textwidth]{media/gan_compress2.png}
\end{center}

\begin{table}[H]
    \begin{center}
    \begin{tabular}{|c|c|c|c|c|c|c|}
    \hline
    \textbf{Image}	& \textbf{ViT-Score} & \textbf{SSIM} & \textbf{MSE} & \textbf{PSNR} & \textbf{FID} & \textbf{BRISQUE} \\
    \hline
    \textbf{Bevo}	 & 0.14 & 0.26 & 3,479.34  &  12.68  &  1,241,999.901  &  39.5535\\
    \textbf{Kliment} & 0.54 & 0.56 & 1,759.28 & 15.64  &  549,089.491 & 44.3570\\
    \textbf{Logo}	 & 0.29 & 0.91 & 393.69 & 22.15 &  81,105.162 & 97.1844\\
    \textbf{Tower}   & 0.03 & 0.53 & 1,315.75 & 16.90 &  331,171.556 & 52.8363\\\hline
    
    \end{tabular}
    \caption[Summary of Results]{ViT-Scores demonstrate a somewhat expected quality assessment.}
    \end{center}
    \end{table}

%==
%	CHAPTER 4
%==
% \cleardoublepage
\chapter{Discussion}

Increase in relevant recent publications (13 alone in 2022 thus far) supports the vision of tying 
vision transformers to image compression.

\ThinHRule

\newpage
\section{Results and Improvements}
A review of the merits from this work.

\subsection{Results}

Features are in deep layers of the GAN network. 
Latent space is hard to decipher
ViT score definition: how many of the top100 labels match
Can take into account probability of label (included in ViT).
While label probability is stable when working with corrupted images (demo), 
unstable when working with generated images.

\subsection{Potential Improvements to Architecture}

Analyze latent space vector with transformer model (not a ViT, but a transformer adaptation)
Steer the GAN faster into training to compress
Slow computational times
Need a GAN trained on all images, not just ImageNet or Celeb 
Need ViT trained on all images
Natural limit to capacity of this model comes from training sets.

Unique positional encoding can also be achieved using trigonometric representation.
For example, a full sentence from text or perhaps a row of pixels from an image could be represented by the various
periods of a sinusoid. Thus, the exact location of each token would be unique.

\section{Optimization}

Experiment and change loss functions (MSE was used, can use a GAN specific loss like FID)

Most valuable technique:
(reducing learning\_rate as the model trains)
changing input images to cater to what the generative model is trained on.
SGD optimizer for GAN

A lot of options still not figured out.
regularization during training: residual dropout, label smoothing

%\newpage

\section{Present and Future of Image Transformers}
Status quo of Transformers in Image Processing, Compression, Analysis, and Generation
Coveted Deep learning based Image compression 
In the deep learning/AI evolutionary process, still too early. Models have not been trained on enough image data.

GAN model only trained on finite set (ImageNet, CIFAR-10, Celeb-HQ faces etc) and resolution.
Need to train GAN on all images ever.

"TransGAN: Two Pure Transformers Can Make One Strong GAN" \citep{jiang2021transgan}
NeurIPS 2021

Goal is to replace Generator and Discriminator in a GAN with Transformers free of 
convolutions. 
Deterministic, Probabilistic
GANs
JPEG/MPEG

\section{Training and Cost Estimates}

Need a ViT on all internet to cost 100M
A 512-core TPU v3 pod costs \$384/hr to use commercially on GCP. 
2.5k core-days means training the ViT cost 24hrs * \$384 * 5 of them = \$46k. 
That's for one of many ViT flavors.
To train a ViT on the whole TACC Frontera at 20k teraflops (top10) or 
Stampede at 10k tflops (top25), it would take respectively about a minute and 2 minutes.
tpu v3 is 420 teraflops * 2500 = 1M Tflops

(like GPT-3 trained on all internet text, Vision T trained on all google images)

(\$100M+), GPT-3 cost \$10M-\$20M
r. In terms of image reconstruction on the decoder side, our image reconstructor needs more
FLOPs than mbt-m due to the feature aggregation module. In
terms of image classification on the decoder side, our image
classifier directly performs inference from the compressed
features without the image reconstruction process, and thus
needs far less computational cost compared with the inference from reconstructed RGB images.


%==
%	CHAPTER 5
%==
% \cleardoublepage
\chapter{Summary}

This chapter serves to conclude this thesis. 

It provides a summary 
of original contributions made by the author while studying and experimenting 
with Vision Transformers (ViT) and Neural Image Compression, as well as the broader 
scientific domains of Machine and Deep Learning and Digital Image Processing. 

A summary of key takeaways is provided for the audience.

The author concludes by acknowledging key contributors to the project and serves closing
remarks.


\ThinHRule

\newpage
\section{Key Contributions}

(experimentation with Vision Transformers, nascent field)
The main merit of this thesis was introducing the ViT-Score.
A Vision Transformer-Assisted metric for evaluating the performance of a neural image compression, the score
correlated with the Generative Adversarial Network verifiably generating a comprehensible image by visual inspection.

The ViT-Score conclusively contributed insights to image quality, relating to human perception of the generated image. 
The metric can also provide additional insights to understanding the latent space (contextual) preservation of an input image.

\NOTE{This work can be viewed as a stepping stone towards an end-to-end Transformer-based
image compression and regeneration.}

\section{Summary}

Throughout this report, the reader was presented with all relevant background knowledge necessary 
to grasp the key contributions listed. 
A Vision Transformer (ViT) was used to evaluate the capacity of 
a GAN to compress and generate an image of choice based on object-level similarities with the original input image.

The new metric, referred to as a ViT-Score, was able to capture and assess the quality of the output images and provide 
valuable insights. The ViT-Score performed well, comparing in capacity to established image quality metrics such as
SSIM, MSE, and PSNR. 


\section{Takeaways}

The future of image compression technology will be based on a Deep Learning methodology.
Due to their generalizability, excellent performance in 1-dimensional data (text),
and proven ability to scale to 2-dimensions, Transformers are an excellent choice of 
architecture to use in image compression and reconstruction.


A Vision Transformer (ViT)-Assisted metric related to image compression can provide 
additional insights to the latent space (contextual) preservation.
Thus, this work can be viewed as a stepping stone towards an end-to-end Transformer-based
image compression and regeneration.


Perhaps, such a metric could be used as a Loss function, embedded within the architecture,
along with being a useful evaluation metric.


It may cost on the order of \$100M and several years, but such a technology will be achieved.


Finally, this image technology can be extended to videos and video compression in further developments.


\section{Acknowledgments}

The author would like to express gratitude towards several individuals and organizations
from The University of Texas at Austin campus.


The major inspiration for this project was gathered from two courses taught by the
reviewers of this thesis.


EE 371Q, Digital Image Processing taught by Professor Alan C. Bovik was the class
where the author learned about 
Image Compression, Image Quality Assessment, and completed a term
project on Generative Adversarial Networks.


CSE 382, Foundations of Machine Learning taught by Professor Rachel A. Ward 
was the class where the author learned key concepts used throughout this thesis
and completed a term project on Vision Transformers (ViT).


Further acknowledgments are made to the Laboratory for Image and Video Engineering (LIVE) at the 
University of Texas at Austin for providing a source for project inspiration and insights.


Finally, the author would like to express gratitude to The Texas Advanced Computing Center (TACC).
TACC provided free access to advanced High-Performance Computing (HPC) resources,
which were used throughout the experimentation process in this thesis.

\section{Closing Remarks}

This thesis is written as a graduation requirement for the degree of Master of Science in

Computational Science, Engineering, and Mathematics awarded by the Oden Institute at 
The University of Texas at Austin.


\NOTE{All code has been made available as open source to the general public in the form of a GitHub repository.}

%==
%	BIBLIOGRAPHY
%==

% \cleardoublepage
\phantomsection %hyperref package support
\addcontentsline{toc}{chapter}{Bibliography} % add entry to table of contents
\pagestyle{plain}


%{\textbf{\LARGE{Bibliography}}}\\	%headline
%\nocitep{*} % Show all Bib-entries (DEBUG)

\bibliographystyle{abbrvnat}
% \bibliography{bib/algorithm.bib,bib/resDes.bib} % for a better structure you can split your bib items into seperate files
\bibliography{bib/resDes.bib}

% \cleardoublepage
\newpage
\phantomsection %hyperref package support
\addcontentsline{toc}{chapter}{Appendix} % add entry to table of contents
\pagestyle{plain}
{\textbf{\LARGE{APPENDIX}}}\\	%headline
\section{Technologies used}


The package \NOTE{algorithm2e} allows you write nearly all imaginable code structures as pseudo code in LaTex. A referece of commands can be found here: \url{http://www.cs.toronto.edu/~frank/Useful/algorithm2e.pdf}. Do not forget to wrap the pseudo code into a figure, thus it will be included into list of figures.\\

The following example demonstrates the package with a bubble sort algorithm. Note, that the directive \NOTE{repeat - until} is the same as the \NOTE{do-while} directive.


\newpage
\subsection{GPU, Local machine}

Writing mathematical equations can be a very important part of your work. This applies especially, if you want to analyze and evaluate your gathered data. The following example shows a calculation of Kohen's Cappa, a measurement which represents the inter coder reliability in coding qualitative data.\\

Let's assume the following results of coding a text by two different persons:


Continuing from this basis, the calculation can be proceeded with Cohen's Kappa defined as


\subsection{TACC, Stampede2, job submission process}

Sometimes you want to include text which include characters that could trigger commands. At this point it useful to wrap them into the \NOTE{verbatim} environment. The text is uninterpreted and listet as is.

\scriptsize
\begin{verbatim}
2015_5_13/23-19-33:550: ************* LOG INFO  *************
2015_5_13/23-19-33:550: name: user1
2015_5_13/23-19-33:550: [0]: 534
2015_5_13/23-19-33:550: [1]: 321
2015_5_13/23-19-33:550: [2]: 094832
2015_5_13/23-19-33:550: [3]: 3980429804
\end{verbatim}
\normalsize	

code


\small
\begin{lstlisting}[keywordstyle=\color{blue},language=Python]
<!-- an XML comment -->
<entry type="normal" pattern="a" />
<entry type="normal" pattern="b" />
<entry type="formula" pattern="t=a+b+c" required="1"/>
\end{lstlisting}
\normalsize	

\end{document}